\subsection{Aditivos}

No contexto deste trabalho, denomina-se aditivo todo material que pode ser incorporado de forma opcional ao procedimento experimental, com o objetivo de ajustar tanto o comportamento da reação quanto as propriedades do sabão. Entre os aditivos descritos a seguir, o etanol e o ácido esteárico foram efetivamente empregados nas práticas, enquanto a casca de laranja, o sal e o açúcar constituem opções comuns que também poderiam ser incorporadas à formulação.

\subsubsection{Etanol}

O etanol atua principalmente como cossolvente na etapa reacional, reduzindo a diferença de polaridade entre a fase oleosa e a fase aquosa alcalina e, com isso, facilitando o contato entre os triglicerídeos e os íons hidróxido. Esse melhor contato entre as fases torna a saponificação mais rápida e homogênea, diminuindo a formação de grumos e de regiões com óleo não reagido. Além disso, o etanol reduz temporariamente a viscosidade da mistura, o que melhora a agitação. Por ser volátil, é removido por aquecimento ao longo do processo; um excesso remanescente ao final pode alterar a textura da barra e prejudicar uma eventual etapa de separação por sal (\textit{salting-out}) \cite{mabrouk2005soap}.

\subsubsection{Ácido esteárico}

O ácido esteárico atua como agente estruturante e de ajuste da formulação. Por ser um ácido graxo livre, reage diretamente com o NaOH formando estearato de sódio, um sal de ácido graxo de cadeia longa e saturada que aumenta a dureza e a consistência do sabão. Esse efeito decorre do fato de cadeias saturadas e longas se empacotarem melhor na matriz sólida, resultando em uma barra mais firme, com menor deformação e maior resistência ao desgaste durante o uso. Em formulações com tendência a um sabão muito mole (como as de alto teor de óleo insaturado), o ácido esteárico ajuda a compensar essa maciez e a melhorar a estabilidade física do produto \cite{alum2024saponification}.

\subsubsection{Casca de laranja}

A casca de laranja, preferencialmente seca e moída, funciona como um aditivo sensorial e funcional leve, podendo conferir aroma cítrico natural, alterar o aspecto visual e proporcionar uma suave ação esfoliante. Essa atuação deve-se à presença de partículas sólidas finas, responsáveis por um efeito abrasivo brando, e de compostos aromáticos naturais, que permanecem parcialmente retidos na matriz do sabão. Por se tratar de material orgânico, deve ser empregada em baixa proporção e com baixa umidade, a fim de reduzir o risco de instabilidade, como o surgimento de odor desagradável, a alteração de cor ou a contaminação \cite{vasave2025herbal}.

\subsubsection{Sal}

O cloreto de sódio (NaCl) é utilizado para o ajuste de dureza e o controle da separação de fases; em baixas concentrações no produto final, tende a deixar a barra mais firme. Quimicamente, o sal reduz a solubilidade dos sais de ácidos graxos em água — o chamado efeito \textit{salting-out} —, favorecendo a separação do sabão da fase aquosa rica em glicerol e sais \cite{shreve1956chemical}; em excesso, porém, pode reduzir a formação de espuma. Por esse motivo, sua dosagem deve ser controlada: uma quantidade muito pequena pode não gerar efeito perceptível, enquanto o excesso pode tornar a barra quebradiça e com menor formação de espuma.

\subsubsection{Açúcar}

O açúcar atua como modificador de espuma e de aparência, em geral aumentando o volume e a persistência da espuma e podendo conferir um aspecto mais translúcido em certas formulações. Esse comportamento está associado ao aumento da fração solúvel na fase aquosa e à modulação da viscosidade local durante a formação do filme de espuma, o que contribui para a estabilização das bolhas. Em excesso, contudo, pode tornar o sabão mais higroscópico e mais mole, de modo que deve ser empregado em faixas moderadas para equilibrar espuma, dureza e estabilidade \cite{mabrouk2005soap}.
